\begin{onehalfspace}

Questa appendice raccoglie il materiale di dettaglio del Capitolo~\ref{chap:obiettivi}: la motivazione crittografica estesa, i requisiti funzionali del sistema sperimentale e la specifica completa dei parametri di input e di output.

\subsection{Motivazione: la Vulnerabilità della Crittografia Attuale}
\label{sec:motivazione}

La quasi totalità delle comunicazioni digitali sicure --- dalle transazioni bancarie alle connessioni \ac{HTTPS}, dai protocolli \ac{SSH} ai certificati digitali --- è protetta da sistemi crittografici la cui sicurezza si fonda su un presupposto computazionale: la difficoltà di fattorizzare il prodotto di due grandi numeri primi. Il sistema \ac{RSA}, ad esempio, costruisce la propria chiave pubblica come $N = p \cdot q$, dove $p$ e $q$ sono primi di centinaia o migliaia di bit. Finché non esiste un algoritmo classico efficiente per recuperare $p$ e $q$ a partire da $N$, la chiave privata rimane al sicuro.

Per comprendere perché i sistemi crittografici attuali siano strutturalmente vulnerabili all'algoritmo di Shor, è necessario mettere a confronto le complessità computazionali in gioco. Il miglior algoritmo classico per la fattorizzazione, il \ac{GNFS}, richiede un tempo:
\begin{equation}
    T_{\text{classico}}(N) = O\!\left(\exp\!\left((c+o(1))(\log N)^{1/3}(\log\log N)^{2/3}\right)\right)
    \label{eq:gnfs_obiettivi}
\end{equation}
che, pur crescendo più lentamente di un esponenziale puro, rimane del tutto impraticabile per i valori di $N$ usati in crittografia. A titolo di esempio, fattorizzare un numero RSA a 2048 bit con algoritmi classici richiederebbe un tempo stimato superiore all'età dell'universo, anche impiegando l'intera potenza di calcolo disponibile sul pianeta.

Nel 1994, Peter Shor dimostrò che un computer quantistico è in grado di eseguire la stessa operazione in tempo \textbf{polinomiale}~\cite{shor1994}:
\begin{equation}
    T_{\text{Shor}}(N) = O\!\left((\log N)^3\right)
    \label{eq:shor_obiettivi}
\end{equation}
Lo scarto tra le due espressioni non è una differenza di grado, ma una differenza di \textit{natura}: il tempo classico cresce in modo super-polinomiale con la lunghezza della chiave, mentre il tempo quantistico cresce polinomialmente. \textbf{Aumentare la lunghezza della chiave RSA non è quindi una contromisura strutturale} contro Shor: raddoppiare $n$ porta il costo asintotico textbook da $O(n^3)$ a $8\,O(n^3)$. Sul piano concreto questo fattore può richiedere molte più risorse fault-tolerant e ritardare un attacco, ma non ripristina un problema computazionalmente difficile per un avversario quantistico scalabile. La stessa vulnerabilità asintotica si estende ai sistemi basati sul logaritmo discreto (Diffie--Hellman, \ac{ECDSA}).

La minaccia ha già una conseguenza operativa: nel modello \textit{harvest now, decrypt later}, dati cifrati oggi e destinati a restare sensibili vengono conservati in vista di una futura macchina crittograficamente rilevante. Non esiste però una previsione affidabile della sua data di arrivo. Le roadmap dei produttori descrivono obiettivi ingegneristici e il \ac{NIST} sottolinea che le stime spaziano da pochi anni a decenni; proprio l'incertezza, unita ai tempi lunghi di migrazione, motiva l'adozione anticipata della \ac{PQC}~\cite{nist_pqc_explainer2026,ibm_roadmap2025}.

Gli ostacoli pratici sono la scala, la profondità fault-tolerant e il \textbf{rumore}. I processori attuali operano in regime \ac{NISQ}: dispongono di molti qubit fisici, ma decoerenza, errori di gate, readout, connettività e overhead di correzione impediscono di sostenere il numero di operazioni logiche richiesto da Shor su chiavi crittografiche. Comprendere questi limiti e stabilire quali strategie li riducano --- e in quale punto della pipeline i modelli di machine learning contribuiscano davvero --- è la ragione d'essere del presente lavoro.

\subsection{Requisiti Funzionali}
\label{sec:requisiti}

Definiti gli obiettivi e il contributo originale, è possibile formalizzare cosa il sistema sperimentale deve concretamente fare. I requisiti seguenti traducono le ipotesi di ricerca in vincoli tecnici verificabili, e costituiscono il riferimento rispetto al quale l'implementazione descritta nel Capitolo~\ref{chap:sviluppo} verrà valutata.

Il sistema sperimentale deve soddisfare i seguenti requisiti funzionali:

\begin{enumerate}

    \item \textbf{Esecuzione configurabile entro il perimetro validato.}
    La campagna rumorosa e la demo pubblica usano l'istanza canonica $N=15$, $a=7$ e un numero configurabile di shot entro limiti dichiarati. Le implementazioni Beauregard per $N=21$ e $N=35$ sono sottoposte a test ideali e audit di risorse, ma non sono esposte come campagne rumorose equivalenti.

    \item \textbf{Modello di rumore parametrizzabile.}
    Il sistema permette di configurare separatamente i parametri Aer $\lambda_{1q}$ e $\lambda_{2q}$, i tempi $T_1/T_2$, il readout asimmetrico della demo e l'eventuale sovrarotazione coerente. I preset della campagna restano uniformi e illustrativi; non riproducono una calibrazione hardware nominata.

    \item \textbf{Raccolta e archiviazione degli output.}
    Il sistema registra la distribuzione di misura del registro di controllo per ogni esecuzione e la archivia in formato strutturato. La riproducibilità usa un seed base e calendari derivati separati per replica, iterazione, simulatore e risoluzione dei pareggi.

    \item \textbf{Analisi statistica del Metodo 1.}
    Il sistema implementa la pipeline del Metodo 1: selezione della misura più frequente dell'istogramma (strategia TOP-1, senza sogliatura preliminare), stima del periodo $r$ tramite frazioni continue, verifica della correttezza dei fattori estratti. Misura e registra il numero di iterazioni necessarie per raggiungere il primo risultato corretto.

    \item \textbf{Training, validazione e inferenza del Metodo 2.}
    Il sistema genera il dataset dal simulatore rumoroso, addestra il classificatore con divisione training/validation/test e ne valuta le metriche. Il modello corrente apprende l'etichetta TOP-1: decide se il candidato dominante conduce ai fattori; se accetta l'istogramma, M2 applica poi la ricerca TOP-4. L'audit TOP-16 è a classe singola e non produce un classificatore.

    \item \textbf{Confronto sistematico sui use case previsti.}
    M1, TOP-4 e M2 vengono eseguiti sugli stessi istogrammi nei due preset $N=15$, producendo metriche appaiate secondo la Sezione~\ref{sec:metriche}. Per $N=21$ e $N=35$ si riportano correttezza ideale e costo compilato; non si afferma che la probabilità rumorosa sia zero, perché la campagna rumorosa non viene eseguita.

    \item \textbf{Riproducibilità.}
    Tutti gli esperimenti sono riproducibili: seed fisso, parametri esplicitati, codice disponibile nel repository del progetto. Questo requisito è condizione necessaria per la validazione accademica dei risultati \cite{shor_hpec2023}.

\end{enumerate}

% -----------------------------------------------
\subsection{Parametri di Input e Output}
\label{sec:input_output}

\subsubsection{Parametri di Input}

La Tabella~\ref{tab:input_params} elenca i parametri di configurazione del sistema. I loro valori vengono fissati per ciascun use case prima dell'esecuzione e rimangono costanti per tutta la durata dell'esperimento corrispondente.

\begin{table}[H]
\centering
\begin{tabular}{@{}lll@{}}
\toprule
\textbf{Parametro}  & \textbf{Unità} & \textbf{Descrizione} \\
\midrule
$N$                 & intero         & Numero da fattorizzare \\
$a$                 & intero         & Base della moltiplicazione modulare ($\gcd(a,N)=1$) \\
$n_{\text{count}}$  & intero         & Qubit del registro di controllo \\
$M_{\text{shots}}$  & intero         & Shot per esecuzione del simulatore \\
$\lambda_{1q}$      & adimensionale  & Parametro Aer del canale depolarizzante 1Q \\
$\lambda_{2q}$      & adimensionale  & Parametro Aer del canale depolarizzante 2Q \\
$T_1$               & ns             & Tempo di rilassamento (amplitude damping) \\
$T_2$               & ns             & Tempo di dephasing \\
$p_{\text{ro}}$     & adimensionale  & Probabilità di readout error \\
\bottomrule
\end{tabular}
\caption{Parametri principali del sistema sperimentale. $\lambda$ segue la convenzione Aer e non coincide con la probabilità aggregata di Pauli non identità.}
\label{tab:input_params}
\end{table}

\subsubsection{Output del Sistema}

Per ogni use case e per ciascuno dei due metodi, il sistema produce le seguenti quantità:

\begin{itemize}
    \item I fattori trovati $p$ e $q$ tali che $p \cdot q = N$, con verifica esplicita della correttezza ($p \neq 1$, $q \neq 1$, $p \cdot q = N$);
    \item Il numero di iterazioni necessarie per il primo risultato corretto;
    \item La distribuzione di frequenza completa degli output del registro di controllo su $M_{\text{shots}}$ esecuzioni;
    \item Le metriche di confronto tra i due metodi definite nella Sezione~\ref{sec:metriche}.
\end{itemize}

Per il solo Metodo 2, vengono prodotte in aggiunta:
\begin{itemize}
    \item L'accuratezza del classificatore e il punteggio F1 sul test set;
    \item La curva ROC con l'area sottostante (\ac{AUC}).
\end{itemize}

\end{onehalfspace}
